\section{Calculation of Antiderivatives}

As previously mentioned, we will use the notation $\int f$ to denote the
antiderivatives of the function $f$. Instead of writing $F\in \int f$ to
indicate that $F$ is an antiderivative of $f$, we may write more simply, albeit
with an abuse of notation, $\int f = F$, remembering (as we did with the
notation of the little $o$) that this relation is not symmetric at all.
Traditionally, one writes
\begin{equation}\label{eq:4137878}
  \int \cos x\, dx = \sin x + c 
\end{equation}
but we could also write more simply:
\mynote{%
The way of writing \eqref{eq:4137878} is only apparently better
than~\eqref{eq:489417} because in some cases not all antiderivatives can be
determined up to a constant (see example~\ref{ex:primitive_non_connesso}).
Therefore, the rule (generally valid for all notations that require
interpretation) is to use the notation of indefinite integrals with caution and
awareness: one must be aware of the meaning of what is written and use a
notation that is understandable to the audience.
}%
\begin{equation}\label{eq:489417}
  \int \cos x\, dx = \sin x.
\end{equation}

\begin{theorem}[Antiderivatives of Some Elementary Functions]
\index{integral!elementary functions}
For every $\alpha \in \RR$, $\alpha\neq -1$, we have
\begin{gather*}
\int x^\alpha\, dx \ni \frac{x^{\alpha+1}}{\alpha+1},
\qquad
\int \frac{1}{x}\, dx \ni \ln\abs{x}
\\
\int e^x \, dx \ni e^x,
\qquad
\int \cos x\, dx \ni \sin x,
\qquad
\int \sin x\, dx \ni -\cos x
\\
\int \cosh x \, dx \ni \sinh x,\qquad
\int \sinh x \, dx \ni \cosh x, \\
\int \frac{1}{1+x^2}\, dx \ni \arctg x, \qquad
\int \frac{1}{\sqrt{1-x^2}}\, dx \ni \arcsin x, \\
\int \frac{1}{\sqrt{x^2+1}}\, dx \ni \settsinh x = \ln\enclose{x+\sqrt{x^2+1}}\\
\int \frac{1}{\sqrt{x^2-1}}\, dx \ni \settcosh x = \ln\enclose{x+\sqrt{x^2-1}}.
\end{gather*}
\end{theorem}
%
\begin{proof}
It suffices to refer to the corresponding table of derivatives of elementary
functions (theorem~\ref{th:derivate_elementari}).
\end{proof}

\begin{theorem}[Linearity of the Indefinite Integral]
For every $\lambda,\mu \in \RR$ and if $f$, $g$ are arbitrary functions, we
have:
\[
  \int \enclose{\lambda f + \mu g} \supset \lambda \int f + \mu \int g
\]
\end{theorem}
%
\begin{proof}
Every element of the set on the right-hand side can be written in the form
$\lambda F + \mu G$ with $F\in \int f$ and $G\in \int g$. Thus, we have $F'=f$
and $G'=g$, from which it follows that
\[
  (\lambda F + \mu G)' = \lambda f + \mu g
\]
and therefore
\[
  \lambda F + \mu G \in \int (\lambda f + \mu g)
\]
as we needed to prove.
\end{proof}

\begin{example}\label{ex:234045}
  Calculate $\int (2x-1)^2\, dx$.
\end{example}
\begin{proof}[Solution.]
  We have 
  \[
  \int (x-1)^2\, dx = \int \enclose{4x^2 - 4x +1}\, dx 
  \ni 4\frac{x^3}{3} - 4 \frac {x^2}{2} + x
  = \frac 4 3 x^3 - 2 x^2 + x.
  \]
\end{proof}

\begin{theorem}[Change of Variable in Integrals]
The following properties hold:
\begin{enumerate}
\item
If $g\colon A \to \RR$ is differentiable and $f\colon g(A) \to \RR$, then
\mymargin{direct substitution}%
\index{substitution!direct}%
\index{integral!direct substitution}%
\begin{equation}\label{eq:sostituzione_1}
  \int f(g(x)) g'(x)\, dx \supset
  \Enclose{\int f(y) \, dy}_{y=g(x)}
\end{equation}
where it is understood that
\[
 \Enclose{F(y)}_{y=g(x)} = F(g(x));
\]

\item
If $g\in C^1([a,b])$ and $f\in C^0(g([a,b]))$, then
\mymargin{change of variable}%
\index{change!of variable}%
\index{integral!change of variable}%
\begin{equation}\label{eq:sostituzione_2}
 \int_{g(a)}^{g(b)} f(x)\, dx = \int_a^b f(g(t))\, g'(t)\, dt;
\end{equation}

\item
If $I$ is an interval, $g\in C^1(I)$ is injective, and $f\in C^0(g(I))$, then
$g^{-1}$ is defined on the interval $g(I)$ and we have
\mymargin{inverse substitution}%
\index{substitution!inverse}%
\index{integral!inverse substitution}%
\begin{equation}\label{eq:239455}
  \int f(x)\, dx \supset \Enclose{\int f(g(t)) g'(t)\, dt}_{t=g^{-1}(x)}
\end{equation}
\end{enumerate}
\end{theorem}
%
\begin{proof}
For~\eqref{eq:sostituzione_1}, it suffices to observe that if $F(y)$ is any
antiderivative of $f(y)$, then $F(g(x))$ is an antiderivative of $f(g(x))\cdot
g'(x)$ (just apply the chain rule for differentiation). Thus, every element of
the set on the right in~\eqref{eq:sostituzione_1} is an element of the set on
the left.

To prove~\eqref{eq:sostituzione_2}, we know that $f$, being continuous, admits
at least one antiderivative $F(x)$. From the previous point, we know that
$F(g(t))$ is an antiderivative of $f(g(t))g'(t)$ (just differentiate to verify
it). Thus, using the fundamental theorem of calculus, we obtain:
\[
\int_{g(a)}^{g(b)}\!\! f(x) \, dx
= \Enclose{F(x)}_{g(a)}^{g(b)}
= F(g(b)) - F(g(a))
= \Enclose{F(g(t))}_a^b
=\int_a^b\! f(g(t))g'(t)\, dt.
\]

To prove~\eqref{eq:239455}, let $F$ be any function element of the set on the
right-hand side of the relation. We will have $F(x) = H(g^{-1}(x))$ with $H(t)$
an antiderivative of $f(g(t))g'(t)$. But then, by the fundamental theorem of
calculus, fixed $t_0 \in I$, we have, by~\eqref{eq:sostituzione_2},
\mynote{If $g^{-1}$ were differentiable, the proof would be immediate, but in general, it may happen that the derivative of $g$ vanishes at some point, and thus $g^{-1}$ might not be differentiable.}%
\begin{align*}
  F(x) - F(g(t_0))
  &= H(g^{-1}(x)) - H(t_0)
  = \int_{t_0}^{g^{-1}(x)} f(g(t)) g'(t)\, dt
  = \int_{g(t_0)}^x f(s)\, ds
\end{align*}
from which, by theorem~\ref{th:torricelli-barrow}, we can differentiate both
sides to obtain $F'(x) = f(x)$. Thus, the function $F$ is an element of the set
on the left-hand side of the relation~\eqref{eq:239455}, as we wanted to prove.
\end{proof}

The formulas of the previous theorem are usually written in the form
\[
  \int f(g(x)) g'(x)\, dx = \int f(y) \, dy
\]
where it is understood that the variables $x$ and $y$ must satisfy the relation
$y=g(x)$ (or, conversely, $x=g^{-1}(y)$). To memorize this formula, it is common
to define the \emph{differential} of a function $g$ as $dg(x) = g'(x)\, dx$
(consistently with the notation $g' = dg / dx$) so that if $y=g(x)$, we have $dy
= g'(x)\, dx$. We will not provide a formal definition of what a differential is
here, but we will certainly use this convenient notation, thinking of it simply
as a typographical convenience.

\begin{example}
Calculate $\int (2x-1)^2\, dx$.
\end{example}
\begin{proof}[Solution.]
We can make the change of variable $y=2x-1$, $dy = 2\, dx$. Thus,
\begin{align*}
\int (2x-1)^2\, dx 
&= \frac 1 2 \int (2x-1)^2\, 2dx
\supset \frac 1 2 \Enclose{\int y^2\, dy}_{y=2x-1}
\ni \frac 1 2 \Enclose{\frac{y^3}{3}}_{y=2x-1}\\
&= \frac 1 6 \enclose{2x-1}^3 = \frac{8x^3 - 12x^3 + 6 x - 1}{6}
= \frac 4 3 x^3 - \frac 2 x^2 + x - \frac 1 6.
\end{align*}
We observe that this antiderivative is different from the one found in
example~\ref{ex:234045}, but it differs from it by a constant. Therefore, both
results are correct.
\end{proof}
\begin{example}
Calculate
$\int \cos^2 x \, dx$.
\end{example}
\begin{proof}[Solution.]
Recalling that $\cos(2t) = \cos^2 t - \sin^2 t = 2\cos^2 t - 1$, we have $\cos^2
t = \frac{1+\cos(2t)}2$ (bisection formula). Thus,
\[
\int \cos^2 t\, dt
= \int\frac{1+\cos(2t)}{2}\, dt
= \int \frac 1 2 \, dt + \int \frac{\cos(2t)}{2}\, dt.
\]
Clearly, $\int \frac 1 2 \, dt  = \frac t 2$.
In the second integral, we can make a change of variable, setting $2t=s$, from
which $2dt = ds$:
\begin{align*}
\int \frac{\cos(2t)}{2}\, dt
&= \frac 1 4 \int \cos(2t)\, 2dt
= \frac 1 4 \Enclose{\int \cos s \, ds}_{s=2t}
= \frac 1 4 \Enclose{\sin s}_{s=2t}\\
&= \frac 1 4 \sin(2t)
= \frac{1}{2}\sin t \cdot \cos t.
\end{align*}
Ultimately, we obtain
\[
  \int \cos^2 x \, dx = \frac{t+\sin t \cdot \cos t}{2}.
\]
\end{proof}


\begin{example}
Calculate
$\int \sqrt{1-x^2}\, dx$.
\end{example}
\begin{proof}
The integrand function is defined for $x\in [-1,1]$. We consider making the
substitution $x=\sin t$ with $t\in [-\pi/2, \pi/2]$. We observe that on
$[-\pi/2,\pi/2]$, the function $\sin t$ is differentiable, invertible, and its
inverse is $t = \arcsin x$. Informally, we have
\[
 x= \sin t, \qquad dx = \cos t \, dt
\]
from which we obtain the formula
\[
 \int \sqrt{1-x^2}\, dx = \Enclose{\int \sqrt{1-\sin^2(t)} \cos t\, dt}_{t=\arcsin x}.
\]
Now observe that for $t\in [-\pi/2, \pi/2]$, we have $\sqrt{1-\sin^2(t)}=\cos
t$, and thus the integral becomes
\[
 \int \sqrt{1-\sin^2(t)}\cos t\, dt = \int \cos^2(t)\, dt.
\]
We calculated this last integral in the previous exercise. Thus, we obtain:
\begin{align*}
 \int \sqrt{1-x^2}\, dx
 &\stackrel{(x=\sin t)}= \int \cos^2(t)\, dt
 = \frac{t + \sin t \cos t}{2} \\
 &= \frac{t + (\sin t) \sqrt{1-\sin^2 t}}{2} \\
 &\stackrel{(t=\arcsin x)}= \frac{\arcsin x + x \sqrt{1-x^2}}{2}.
\end{align*}
\end{proof}

Note that the graph $y=\sqrt{1-x^2}$ from the previous example is a semicircle,
and thus the integral
\[
 \int_{-1}^1 \sqrt{1-x^2}\, dx = \Enclose{\frac{\arcsin x +x\sqrt{1-x^2}}{2}}_{-1}^{1}
 = \frac{\pi}{2}.
\]
represents the area of a unit semicircle.

\begin{exercise}
Calculate
\[
 \int \sqrt{1+x^2}\, dx, \qquad \int \sqrt{x^2-1}\, dx.  
\]

Hint: Analogous to the integral of $\sqrt{1-x^2}$, you can perform the
substitution $t=\sinh x$ in the first integral and $t=\cosh x$ in the second.
\end{exercise}

\begin{theorem}[Integration by Parts]
\mymark{*}%
\mymargin{integral by parts}%
\index{integral!by parts}%
Let $f\colon A\subset \RR \to \RR$ be any function, let $g\colon A \to\RR$ be a
differentiable function, and let $F \in \int f$. Then
\[
  \int f\cdot g \supset F \cdot g - \int F \cdot g'.
\]

In particular, if $f\in C^0([a,b])$ and $g\in C^1([a,b])$ and $F \in \int f$, we
have
\[
  \int_a^b f\cdot g = \Enclose{F\cdot g}_a^b - \int_a^b F \cdot g'.
\]
\end{theorem}
%
\begin{proof}
\mymark{*}
Every function in the set on the right can be written in the form $F\cdot g - H$
with $H \in \int F \cdot g'$. Thus, $H' = F \cdot g'$ and, by hypothesis,
$F'=f$, from which it follows that
\[
(F\cdot g - H)' = F' \cdot g + F \cdot g' - H' = F' \cdot g
= f\cdot g
\]
which is what we needed to prove.

The second part of the theorem follows directly from the fundamental theorem of
calculus (valid since both $f\cdot g$ and $F \cdot g'$ are continuous
functions), observing that
\[
\Enclose{F\cdot g - \int F \cdot g'}_a^b
= \Enclose{F\cdot g}_a^b - \int_a^b F \cdot g'.
\]
\end{proof}

\begin{example}
We want to calculate
\[
  \int x \cos x\, dx.
\]
The method of integration by parts allows us to reduce the integral of a product
to an integral of a product where one of the factors is integrated and the other
is differentiated. In this case, it will be convenient to differentiate the
factor $x$ and integrate the factor $\cos x$ so as to reduce to the integral of
$1\cdot \sin x$, which we know how to solve. Specifically, we have
\[
 \int x \cos x\, dx = x \sin x - \int 1 \cdot \sin x \, dx
  = x \sin x + \cos x.
\]
\end{example}

\begin{example}
  \label{ex:822309}
We want to calculate
\[
 \int e^x \cos x\, dx.
\]
In this case, if we use integration by parts, we can reduce this integral to
$\int e^x \sin x$. Integrating again by parts, we will again reduce to $\int e^x
\cos x$. However, if in these steps we obtain the original quantity with a
changed sign, we can solve the resulting equation to find the desired result.

Specifically:
\begin{align*}
\int e^x \cos x\, dx
&= e^x \sin x - \int e^x \sin x\, dx\\
 &= e^x \sin x - \Enclose{e^x(-\cos x) - \int e^x(-\cos x)\, dx} \\
 &= e^x \sin x + e^x \cos x - \int e^x \cos x \, dx
\end{align*}
from which:
\[
 2 \cdot \int e^x \cos x\, dx  = e^x \sin x + e^x \cos x
\]
that is,
\[
  \int e^x \cos x\, dx = \frac{e^x(\sin x + \cos x)}{2}.
\]
\end{example}

\begin{theorem}[More Antiderivatives of Elementary Functions]
\mymark{***}
We have
\begin{align*}
  \int \ln x\, dx  &= x \ln x - x, \\
  \int \arctg x\, dx &= x \arctg x - \ln \sqrt{1+x^2},\\
  \int \arcsin x\, dx &= x \arcsin x + \sqrt{1-x^2},\\
  \int \arccos x\, dx &= x \arccos x - \sqrt{1-x^2}.
\end{align*}
\end{theorem}
%
\begin{proof}
\mymark{***}
The idea, in all cases, is that the derivative of the integrand function is much
simpler than the function itself (in the first two cases, it is a rational
function, in the other two, an irrational but not transcendental function). It
may therefore be useful to apply integration by parts in the form:
\[
  \int f(x)\, dx = \int 1\cdot f(x)\, dx = x f(x) - \int x f'(x)\, dx.
\]
In the first case, we have:
\[
\int \ln x\, dx = x \ln x - \int x \frac{1}{x}\, dx
 = x \ln x - \int 1 dx = x \ln x - x.
\]
In the second case:
\[
\int \arctg x\, dx = x \arctg x - \int \frac{x}{1+x^2}\, dx.
\]
We then perform a change of variable $y=1+x^2$:
\begin{align*}
\int \frac{x}{1+x^2}\, dx
&= \frac{1}{2}\int \frac{1}{1+x^2} 2x \, dx
\stackrel{y=1+x^2} = \frac{1}{2}\int \frac 1 y\, dy \\
&= \frac{\ln y}{2} = \frac{1}{2}\ln\enclose{1+x^2}
\end{align*}
from which, in conclusion:
\[
 \int \arctg x\, dx = x \arctg x - \frac 1 2 \ln\enclose{1+x^2}.
\]
In the third case:
\[
 \int 1\cdot \arcsin x\, dx = x\cdot \arcsin x - \int\frac{x}{\sqrt{1-x^2}}\, dx  
\]
and by performing the substitution $y=1-x^2$, we arrive at 
\[
  \int \arcsin x\, dx = x \cdot \arcsin x + \sqrt{1-x^2}.  
\]
The last integral is carried out similarly.
\end{proof}

\begin{remark}
If the method of integration by parts is applied in the search for the
antiderivative of the function $\frac{1}{x\ln x}$, integrating the factor $\frac
1 x$ and differentiating the factor $\frac 1 {\ln x}$, a seemingly disconcerting
phenomenon occurs:
\begin{align*}
  \int \frac{1}{x\ln x}\,dx
    &= \ln x \cdot \frac{1}{\ln x} - \int \ln x \cdot \frac{\frac 1 x}{-\ln^2
  x}\, dx\\
  &= 1 + \int \frac{1}{x\ln x}\, dx.
\end{align*}
One might think of simplifying the integrals on both sides of the equality to
obtain $0=1$. This cannot be done because, let us remember, the indefinite
integral is a set of functions (the antiderivatives of the function
$\frac{1}{x\ln x}$ in this case), and we know that adding a constant to an
antiderivative yields another antiderivative. Therefore, it should not be
surprising that adding $1$ to the set of antiderivatives leaves it unchanged.

For the record: the integral in question can be calculated by substitution,
setting $y=\ln x$ and finding $F(x) = \ln\abs{\ln x}$ as one of the
antiderivatives.
\end{remark}